\chapter{Functional, Fourier, and Complex Analysis}

\section{Hilbert Spaces and Square-Integrable Functions}

To rigorously generalize the concept of Euclidean space $\mathbb{R}^N$ to infinite dimensions, we must construct a vector space of functions that retains the geometric notions of length (norm) and angle (inner product), while ensuring that the space is \textit{complete}. Completeness guarantees that limits of Cauchy sequences exist within the space, which is an absolute prerequisite for performing calculus and infinite series expansions. The standard choice for physical applications is the space of square-integrable functions, denoted $L^2(\mathbb{R})$.

However, constructing this space requires addressing several foundational subtleties that are often glossed over. We will build it step-by-step, explicitly justifying every definition.

\subsection{Measure-Theoretic Construction of \texorpdfstring{$L^2(\mathbb{R})$}{L2(R)}}

Before defining the space, we must address the integral. In elementary calculus, the Riemann integral is used. However, the Riemann integral is inadequate for our purposes because the space of Riemann-integrable functions is \textit{not complete} under the $L^2$ norm. Specifically, a sequence of continuous (and thus Riemann-integrable) functions can converge in the $L^2$ norm to a highly discontinuous function that is not Riemann-integrable. To ensure completeness, we must use the \textit{Lebesgue integral}.
\begin{itemize}
\item \textit{Lebesgue measure ($\mu$)}: a function that assigns a ``length'' or ``volume'' to subsets of $\mathbb{R}$. Unlike the intuitive length of an interval, it is defined on a $\sigma$-algebra of sets, allowing us to measure highly fragmented sets (like the rational numbers, which have measure zero).
\item \textit{Measurable functions}: a function $v: \mathbb{R} \to \mathbb{C}$ is Lebesgue measurable if the preimage of any open set in $\mathbb{C}$ is a Lebesgue measurable set in $\mathbb{R}$. This is the exact class of functions for which the Lebesgue integral

\[
\int_\mathbb{R} |v(t)|^2 \dd t
\]
is well-defined.
\end{itemize}


Let $\mathcal{L}^2(\mathbb{R})$ be the set of all Lebesgue measurable functions $v: \mathbb{R} \to \mathbb{C}$ such that the integral of their squared modulus is finite:
\begin{equation}
\mathcal{L}^2(\mathbb{R}) = \left\{ v \in \mathcal{M}(\mathbb{R}) : \int_\mathbb{R} |v(t)|^2 \dd t < \infty \right\}
\end{equation}
where $\mathcal{M}(\mathbb{R})$ denotes the set of all complex-valued Lebesgue measurable functions.

\begin{remark}[The Problem of Positive Definiteness]
If we attempt to define a norm on $\mathcal{L}^2(\mathbb{R})$ as
\[
\|v\| = \left( \int_\mathbb{R} |v(t)|^2 \dd t \right)^{1/2},
\]
we encounter a fatal flaw. The integral of a nonnegative function is zero if and only if the function is zero \textit{almost everywhere} (a.e.), meaning the set of points where it is nonzero has Lebesgue measure zero.
Therefore, $\|v\| = 0$ does \textit{not} imply $v(t) = 0$ for all $t \in \mathbb{R}$; it only implies $v(t) = 0$ except possibly on a set of measure zero. This violates the axiom of positive definiteness for a norm, which strictly requires
\[
\|v\| = 0 \implies v = \mathbf{0}.
\]

\end{remark}

To resolve this, we must introduce an equivalence relation and form a quotient space.

\begin{definition}[Equivalence Relation and Quotient Space]
We define a relation $\sim$ on $\mathcal{L}^2(\mathbb{R})$ such that $v \sim u$ if and only if $v(t) = u(t)$ for almost every $t \in \mathbb{R}$. Formally, the set $\{t \in \mathbb{R} : v(t) \neq u(t)\}$ has Lebesgue measure zero.
This relation is reflexive, symmetric, and transitive, making it a valid equivalence relation. It partitions $\mathcal{L}^2(\mathbb{R})$ into disjoint equivalence classes.
We define the Hilbert space $L^2(\mathbb{R})$ as the quotient space:
\begin{equation}
\begin{aligned}
L^2(\mathbb{R}) &= \mathcal{L}^2(\mathbb{R}) / \\
&\sim
\end{aligned}
\end{equation}

An element of $L^2(\mathbb{R})$ is not a single function, but an equivalence class $[v]$ of functions that differ only on sets of measure zero. By convention, we abuse notation and write $v$ instead of $[v]$, treating functions that are equal almost everywhere as identical elements of the space.
\end{definition}

\subsection{The Inner Product and Norm}

With the elements of the space properly defined as equivalence classes, we can define the geometric structure.

\begin{definition}[Inner Product on $L^2(\mathbb{R})$]
For any two elements $v, u \in L^2(\mathbb{R})$, we define the inner product as:
\begin{equation}
\langle v, u \rangle = \int_\mathbb{R} v(t) \overline{u(t)} \dd t
\end{equation}

\end{definition}

\begin{remark}[Well-definedness and Finiteness of the Inner Product]
Before using this formula, we must prove the integral actually converges to a finite complex number. We apply the \textit{Cauchy--Schwarz inequality for integrals}, which states that for measurable functions $f, g$:
\begin{equation}
\int_\mathbb{R} |f(t) g(t)| \dd t \leq \left( \int_\mathbb{R} |f(t)|^2 \dd t \right)^{1/2} \left( \int_\mathbb{R} |g(t)|^2 \dd t \right)^{1/2}
\end{equation}
Setting $f(t) = v(t)$ and $g(t) = \overline{u(t)}$, we get:
\begin{equation}
\int_\mathbb{R} |v(t) \overline{u(t)}| \dd t \leq \|v\| \|u\| < \infty
\end{equation}
Since the integral of the absolute value is finite, the integral
\[
\int v \bar{u}
\]
is absolutely convergent and thus well-defined. Furthermore, because $v$ and $u$ are equivalence classes, changing the representative functions on a set of measure zero does not change the value of the integral, making the inner product well-defined on the quotient space.
\end{remark}

The inner product induces a norm via
\[
\|v\| = (\langle v, v \rangle)^{1/2}.
\]
Because we quotiented out the null sets,
\[
\|v\| = 0 \implies \int |v|^2 = 0 \implies v = 0
\]
a.e. $\implies v$ is the zero vector in $L^2(\mathbb{R})$. The triangle inequality for this norm follows directly from the \textit{Minkowski inequality}.

\subsection{Completeness and the Riesz--Fischer Theorem}

A normed space is a \textit{Banach space} if it is complete. A \textit{Hilbert space} is a Banach space whose norm is induced by an inner product.

\begin{definition}[Cauchy Sequence and Completeness]
A sequence $\{v_n\}$ in $L^2(\mathbb{R})$ is a \textit{Cauchy sequence} if for every $\varepsilon > 0$, there exists an integer $N$ such that for all $m, n > N$, $\|v_m - v_n\| < \varepsilon$.
The space is \textit{complete} if every Cauchy sequence converges to a limit $v \in L^2(\mathbb{R})$, meaning
\[
\lim_{n \to \infty} \|v_n - v\| = 0.
\]

\end{definition}

\begin{theorem}[Riesz--Fischer Theorem]
The space $L^2(\mathbb{R})$ is complete with respect to the norm induced by its inner product. Consequently, $L^2(\mathbb{R})$ is a Hilbert space.
\end{theorem}

\begin{proof}[Rigorous Sketch of Logic]
The original text's reliance on the Dominated Convergence Theorem (DCT) is a common logical leap; DCT requires a \textit{pre-existing} dominating function, which a general Cauchy sequence does not possess a priori. The rigorous proof constructs the limit function explicitly using a subsequence.

\textit{Step 1: Extract a rapidly converging subsequence.}
Since $\{v_n\}$ is Cauchy, we can find a subsequence $\{v_{n_k}\}$ such that
\[
\|v_{n_{k+1}} - v_{n_k}\| < 2^{-k}
\]
for all $k \geq 1$.

\textit{Step 2: Define a sequence of partial sums of differences.}
Define the partial sums
\[
g_K(t) = \sum_{k=1}^K |v_{n_{k+1}}(t) - v_{n_k}(t)|.
\]
This is a monotonically increasing sequence of nonnegative measurable functions.

\textit{Step 3: Apply the Monotone Convergence Theorem (MCT).}
By Minkowski's inequality, the $L^2$ norm of $g_K$ is bounded:
\begin{equation}
\|g_K\| \leq \sum_{k=1}^K \|v_{n_{k+1}} - v_{n_k}\| < \sum_{k=1}^K 2^{-k} < 1
\end{equation}
Thus,
\[
\int_\mathbb{R} |g_K(t)|^2 \dd t < 1.
\]
By the Monotone Convergence Theorem, the pointwise limit
\[
g(t) = \lim_{K \to \infty} g_K(t)
\]
exists, and
\[
\int_\mathbb{R} |g(t)|^2 \dd t \leq 1.
\]
Because the integral of $|g(t)|^2$ is finite, $g(t)$ must be finite almost everywhere.

\textit{Step 4: Construct the limit function.}
Since $g(t) < \infty$ a.e., the series
\[
\sum_{k=1}^\infty |v_{n_{k+1}}(t) - v_{n_k}(t)|
\]
converges absolutely a.e. This implies that the telescoping series:
\begin{equation}
v_{n_1}(t) + \sum_{k=1}^\infty \left( v_{n_{k+1}}(t) - v_{n_k}(t) \right)
\end{equation}
converges absolutely a.e. to some measurable function $v(t)$. By construction, $v_{n_k}(t) \to v(t)$ pointwise a.e.

\textit{Step 5: Prove $v \in L^2(\mathbb{R})$ and $v_n \to v$ in norm.}
Using Fatou's Lemma on the sequence $|v_{n_k} - v_m|^2$ (for fixed $m$), one can show that $v - v_m \in L^2(\mathbb{R})$, which implies $v \in L^2(\mathbb{R})$. Finally, because $\{v_n\}$ is Cauchy and a subsequence converges to $v$ in norm, the entire sequence must converge to $v$ in the $L^2$ norm.
\end{proof}

\subsection{Separability and Orthonormal Bases}

Completeness allows us to treat $L^2(\mathbb{R})$ as the rigorous infinite-dimensional analogue of $\mathbb{R}^N$. In $\mathbb{R}^N$, any vector is expanded in a finite orthonormal basis. In infinite dimensions, we require a countable basis.

\begin{definition}[Separability]
A metric space is \textit{separable} if it contains a countable, dense subset. $L^2(\mathbb{R})$ is separable because the set of step functions with rational endpoints and rational heights is countable and dense in $L^2(\mathbb{R})$.
\end{definition}

Because $L^2(\mathbb{R})$ is a separable Hilbert space, it admits a \textit{countable orthonormal basis} (more formally, a complete orthonormal system) $\{\phi_n\}_{n \in \mathbb{N}}$. ``Orthonormal'' means $\langle \phi_n, \phi_m \rangle = \delta_{nm}$ (the Kronecker delta). ``Complete'' means the only vector orthogonal to all $\phi_n$ is the zero vector.

Just as in $\mathbb{R}^N$, any $v \in L^2(\mathbb{R})$ has an expansion with coefficients given by its inner products with the basis functions:
\begin{equation}
\begin{aligned}
v &= \sum_{n=1}^\infty c_n \phi_n, \\
c_n &= \langle v, \phi_n \rangle.
\end{aligned}
\end{equation}

\begin{remark}[Nature of Convergence]
The infinite sum
\[
\sum_{n=1}^\infty c_n \phi_n
\]
does \textit{not} necessarily converge pointwise. The equality $v = \sum c_n \phi_n$ is understood in the \textit{$L^2$-norm sense} (mean-square convergence). Let
\[
S_N = \sum_{n=1}^N c_n \phi_n
\]
be the partial sum. The expansion means:
\begin{equation}
\lim_{N \to \infty} \left\| v - S_N \right\| = 0 \implies \lim_{N \to \infty} \int_\mathbb{R} \left| v(t) - \sum_{n=1}^N c_n \phi_n(t) \right|^2 \dd t = 0
\end{equation}

This is guaranteed by the completeness of the space and Bessel's inequality becoming an equality (Parseval's identity):
\[
\|v\|^2 = \sum_{n=1}^\infty |c_n|^2.
\]

\end{remark}

Figure~\ref{fig:atlas-orthogonal-modes} illustrates orthogonal oscillatory modes in a function space.

\begin{figure}[!htbp]
\centering
\begin{tikzpicture}[>=Stealth]
\begin{axis}[
  width=12.5cm,height=6.4cm,
  axis lines=left,
  ticklabel style={font=\small},label style={font=\small},
  xlabel={$x$},ylabel={$e_n(x)$},
  xmin=-0.03,xmax=1.03,ymin=-1.6,ymax=1.65,
  xtick={0,0.5,1},ytick={-1,0,1},
  legend style={at={(0.5,1.02)},anchor=south,legend columns=3,
    font=\scriptsize,draw=black!20,fill=white,column sep=1em}
]
\draw[black!20,densely dashed] (axis cs:0,-1.5) -- (axis cs:0,1.5);
\draw[black!20,densely dashed] (axis cs:1,-1.5) -- (axis cs:1,1.5);
\addplot[figblue,very thick,domain=0:1,samples=140] {sqrt(2)*sin(deg(pi*x))};
\addlegendentry{$n=1$}
\addplot[figred,thick,dashed,domain=0:1,samples=140] {sqrt(2)*sin(deg(2*pi*x))};
\addlegendentry{$n=2$}
\addplot[figgreen,thick,dashdotted,domain=0:1,samples=140] {sqrt(2)*sin(deg(3*pi*x))};
\addlegendentry{$n=3$}
\end{axis}
\end{tikzpicture}
\caption{The first three functions $\sqrt{2}\sin(n\pi x)$ on $[0,1]$, extended by zero outside this interval, are orthonormal in $L^2(\mathbb{R})$. Their different oscillation patterns have zero pairwise inner products. Orthonormality refers to an integral, not to perpendicular curves in the picture.}
\label{fig:atlas-orthogonal-modes}
\end{figure}

\section{Integral Operators and Spectral Theory}

\subsection{Hilbert--Schmidt Integral Operators}

In linear algebra, linear transformations on $\mathbb{R}^N$ are represented by matrices. In $L^2(\mathbb{R})$, a broad and highly useful class of linear operators is represented by integral kernels.

\begin{definition}[Linear Integral Operator]
Let $W(t, t')$ be a measurable function on $\mathbb{R} \times \mathbb{R}$. We define a linear operator
\[
W: L^2(\mathbb{R}) \to L^2(\mathbb{R})
\]
acting on a function $v \in L^2(\mathbb{R})$ via:
\begin{equation}
(Wv)(t) = \int_\mathbb{R} W(t, t') v(t') \dd t'
\end{equation}

\end{definition}

For $W$ to be a valid operator on $L^2(\mathbb{R})$, it must map $L^2$ functions to $L^2$ functions, and ideally, it should be a \textit{bounded operator}. An operator is bounded if there exists a constant $C > 0$ such that $\|Wv\| \leq C\|v\|$ for all $v \in L^2(\mathbb{R})$. Boundedness is equivalent to continuity in normed spaces.

\begin{definition}[Hilbert--Schmidt Kernel and Operator]
If the kernel $W(t, t')$ is square-integrable over the product space $\mathbb{R}^2$ with respect to the product Lebesgue measure $\dd t \dd t'$, meaning:
\begin{equation}
\|W\|_{HS}^2 := \iint_{\mathbb{R}^2} |W(t, t')|^2 \dd t \dd t' < \infty
\end{equation}
then $W$ is called a \textit{Hilbert--Schmidt kernel}.
\end{definition}

\begin{theorem}[Boundedness of Hilbert--Schmidt Operators]
Every Hilbert--Schmidt kernel defines a bounded linear operator from $L^2(\mathbb{R})$ to $L^2(\mathbb{R})$, with operator norm bounded by $\|W\|_{HS}$.
\end{theorem}

\begin{proof}[Explicit Derivation of Boundedness]
We must prove that if $v \in L^2(\mathbb{R})$, then
\[
\begin{gathered}
Wv \in L^2(\mathbb{R}) \\
\|Wv\| \leq \|W\|_{HS} \|v\|.
\end{gathered}
\]

First, we evaluate the pointwise absolute value of $(Wv)(t)$ using the \textit{Cauchy--Schwarz inequality for integrals} with respect to the variable $t'$:
\begin{align}
|(Wv)(t)|^2 &= \left| \int_\mathbb{R} W(t, t') v(t') \dd t' \right|^2 \\
    &\leq \left( \int_\mathbb{R} |W(t, t')|^2 \dd t' \right) \left( \int_\mathbb{R} |v(t')|^2 \dd t' \right) \\
    &= \|v\|^2 \int_\mathbb{R} |W(t, t')|^2 \dd t'
\end{align}
Notice that we factored out $\|v\|^2$ because it does not depend on $t'$. Next, to find the $L^2$ norm of $Wv$, we integrate both sides of this inequality with respect to $t$ over $\mathbb{R}$:
\begin{align}
\|Wv\|^2 = \int_\mathbb{R} |(Wv)(t)|^2 \dd t &\leq \int_\mathbb{R} \left( \|v\|^2 \int_\mathbb{R} |W(t, t')|^2 \dd t' \right) \dd t \\
    &= \|v\|^2 \iint_{\mathbb{R}^2} |W(t, t')|^2 \dd t' \dd t \\
    &= \|v\|^2 \|W\|_{HS}^2
\end{align}
Taking the square root of both sides yields:
\begin{equation}
\|Wv\| \leq \|W\|_{HS} \|v\|
\end{equation}

This explicitly proves that $W$ is a bounded linear operator, and the constant $C$ in the definition of boundedness is precisely the Hilbert--Schmidt norm of the kernel, $\|W\|_{HS}$. Furthermore, because $|Wv(t)|^2$ is bounded by an $L^1$ function of $t$, it guarantees that $Wv \in L^2(\mathbb{R})$.
\end{proof}

Figure~\ref{fig:atlas-integral-kernel} visualizes a symmetric square-integrable kernel on a bounded domain.

\begin{figure}[!htbp]
\centering
\begin{tikzpicture}[>=Stealth]
\begin{axis}[width=8cm,height=7cm,axis lines=middle,ticklabel style={font=\small},label style={font=\small},legend style={font=\scriptsize,draw=black!20,fill=white},view={0}{90},axis lines=box,xlabel={$x$},ylabel={$y$},xtick={0,0.5,1},ytick={0,0.5,1},colorbar,colormap={kernel}{color=(white);color=(figblue)},colorbar style={ylabel={$K(x,y)$},ytick={0,0.5,1}},point meta min=0,point meta max=1]
\addplot3[surf,shader=interp,domain=0:1,y domain=0:1,samples=25,samples y=25] {min(x,y)};
\end{axis}
\end{tikzpicture}
\caption{The color map represents $K(x,y)=\min(x,y)$ on the unit square, with zero extension outside it. Symmetry across the diagonal corresponds to $K(x,y)=K(y,x)$. Each fixed-$x$ slice supplies the weights used to integrate an input function; the bounded, compactly supported kernel is square-integrable.}
\label{fig:atlas-integral-kernel}
\end{figure}

\subsection{Eigenfunctions, Self-Adjointness, and Compactness}

In finite-dimensional linear algebra, the Spectral Theorem states that any real symmetric matrix (or complex Hermitian matrix) is diagonalizable by a unitary matrix. To generalize this to infinite-dimensional Hilbert spaces, we must carefully identify which properties of finite matrices survive the transition to infinite dimensions, and which new conditions are required to prevent pathological behavior.


Let $\calH$ be a separable Hilbert space with inner product $\langle \cdot, \cdot \rangle$ and induced norm $\|\cdot\|$. Let $W: \calH \to \calH$ be a bounded linear operator. We first define the adjoint and then the two crucial properties required for the spectral theorem.

\begin{definition}[Adjoint and Self-Adjoint Operator]
The \textit{adjoint} of $W$, denoted $W^\dagger$, is the unique bounded linear operator satisfying $\langle Wv, u \rangle = \langle v, W^\dagger u \rangle$ for all $v, u \in \calH$.
An operator $W$ is \textit{self-adjoint} (or Hermitian) if $W = W^\dagger$, meaning that the following identity holds for all $v, u \in \calH$:
\begin{equation}
\langle Wv, u \rangle = \langle v, Wu \rangle.
\end{equation}

\end{definition}

\begin{remark}[Why Self-Adjointness?]
Self-adjointness is the infinite-dimensional analogue of matrix symmetry and yields two key spectral consequences:
\textit{1. Real Eigenvalues:} If $Wv = \lambda v$ for $v \neq 0$, then

\[
\lambda \langle v, v \rangle = \langle Wv, v \rangle = \langle v, Wv \rangle = \langle v, \lambda v \rangle = \bar{\lambda} \langle v, v \rangle.
\]
Since $\langle v, v \rangle > 0$, we must have $\lambda = \bar{\lambda}$, so $\lambda \in \R$.
\textit{2. Orthogonal Eigenvectors:} If $Wv_1 = \lambda_1 v_1$ and $Wv_2 = \lambda_2 v_2$ with $\lambda_1 \neq \lambda_2$, then

\[
\lambda_1 \langle v_1, v_2 \rangle = \langle Wv_1, v_2 \rangle = \langle v_1, Wv_2 \rangle = \lambda_2 \langle v_1, v_2 \rangle.
\]
Thus,
\[
(\lambda_1 - \lambda_2)\langle v_1, v_2 \rangle = 0,
\]
which forces $\langle v_1, v_2 \rangle = 0$.
\end{remark}

\begin{definition}[Compact Operator]
An operator $W: \calH \to \calH$ is \textit{compact} if it maps every bounded set in $\calH$ to a relatively compact set (a set whose closure is compact). In a metric space like a Hilbert space, compactness is equivalent to sequential compactness. Therefore, $W$ is compact if and only if for every bounded sequence $\{v_n\} \subset \calH$, the image sequence $\{Wv_n\}$ contains a convergent subsequence in $\calH$.
\end{definition}

\begin{remark}[Why Compactness?]
In finite dimensions, the closed unit ball is compact (Heine--Borel theorem), so every bounded linear operator is automatically compact. In infinite dimensions, the closed unit ball is \textit{never} compact. Compact operators are those that ``compress'' the infinite-dimensional space enough to mimic finite-dimensional behavior. Without compactness, the spectrum of an operator can be continuous, meaning the operator might have no eigenvectors in $\calH$ at all.
\end{remark}

\subsection{The Spectral Theorem for Compact Self-Adjoint Operators}

\begin{theorem}[Spectral Theorem]
Let $W$ be a compact self-adjoint operator on a separable Hilbert space $\calH$. Then:
\begin{enumerate}
\item \textit{Real eigenvalues}: the eigenvalues of $W$ are all real.
\item \textit{Orthonormal eigenvectors}: there exists an orthonormal sequence of eigenvectors $\{\phi_n\}$ with corresponding real eigenvalues $\{\lambda_n\}$.
\item \textit{Eigenvalue convergence}: if the sequence of nonzero eigenvalues is infinite, then

\[
\lim_{n \to \infty} \lambda_n = 0.
\]

\item \textit{Spectral expansion}: any vector $v \in \calH$ can be expanded as:

\begin{equation}
Wv = \sum_{n} \lambda_n \langle v, \phi_n \rangle \phi_n
\end{equation}
where the series converges in the norm of $\calH$. (If $W$ is injective, $\{\phi_n\}$ forms a complete orthonormal basis for $\calH$.)
\end{enumerate}
\end{theorem}

\begin{proof}[Explicit Logic for $\lambda_n \to 0$]
We must explicitly prove why compactness forces the eigenvalues to converge to zero.
Suppose, for contradiction, that $\lambda_n \not\to 0$. Then there exists some $\varepsilon > 0$ and a subsequence of eigenvalues $\{\lambda_{n_k}\}$ such that $|\lambda_{n_k}| \geq \varepsilon$ for all $k$.
Consider the corresponding eigenvectors $\{\phi_{n_k}\}$. Because they are orthonormal, they form a bounded sequence ($\|\phi_{n_k}\| = 1$).
Now, examine the image sequence under $W$: $W\phi_{n_k} = \lambda_{n_k} \phi_{n_k}$.
For any $j \neq k$, the distance between two terms in this image sequence is:
\begin{align}
\|W\phi_{n_j} - W\phi_{n_k}\|^2 &= \|\lambda_{n_j} \phi_{n_j} - \lambda_{n_k} \phi_{n_k}\|^2 \notag \\
    &= \langle \lambda_{n_j} \phi_{n_j} - \lambda_{n_k} \phi_{n_k}, \lambda_{n_j} \phi_{n_j} - \lambda_{n_k} \phi_{n_k} \rangle \notag \\
    &= |\lambda_{n_j}|^2 \|\phi_{n_j}\|^2 + |\lambda_{n_k}|^2 \|\phi_{n_k}\|^2 - 2\Re(\lambda_{n_j} \bar{\lambda}_{n_k} \langle \phi_{n_j}, \phi_{n_k} \rangle)
\end{align}
Because the eigenvectors are orthonormal, $\langle \phi_{n_j}, \phi_{n_k} \rangle = 0$ and $\|\phi_{n_k}\| = 1$. Thus:
\begin{equation}
\begin{aligned}
\|W\phi_{n_j} - W\phi_{n_k}\|^2 &= |\lambda_{n_j}|^2 + |\lambda_{n_k}|^2 \\
&\geq \varepsilon^2 + \varepsilon^2 \\
&= 2\varepsilon^2
\end{aligned}
\end{equation}

This means the distance between any two distinct terms in the sequence $\{W\phi_{n_k}\}$ is at least $\sqrt{2}\varepsilon$. Consequently, this sequence cannot contain any Cauchy subsequence, and therefore cannot contain any convergent subsequence. This directly contradicts the definition of $W$ being a compact operator. Hence, our assumption was false, and we must have $\lambda_n \to 0$.
\end{proof}

\begin{remark}[Connection to Finite-Dimensional Diagonalization]
In finite dimensions, diagonalization is written as $W = E \Lambda E^{-1}$, where $E$ is the matrix of eigenvectors and $\Lambda$ is the diagonal matrix of eigenvalues. Because $W$ is self-adjoint, $E$ is unitary, so $E^{-1} = E^\dagger$.
The infinite-dimensional expansion
\[
Wv = \sum_n \lambda_n \langle v, \phi_n \rangle \phi_n
\]
is the exact analogue. If we define the operator $E$ by its action
\[
E c = \sum_n c_n \phi_n,
\]
then $W = E \Lambda E^\dagger$, where $\Lambda$ acts by multiplying the $n$-th component by $\lambda_n$.
\end{remark}

\subsection{Convolution Operators and the Continuous Spectrum}

When we drop the requirement of compactness, the spectral theorem fails in its discrete form. A prime example in physics and signal processing is the convolution operator.

\begin{definition}[Convolution Operator]
Let $K \in L^1(\R)$ be a kernel function. The convolution operator $W_K: L^2(\R) \to L^2(\R)$ is defined by:
\begin{equation}
\begin{aligned}
(W_K v)(t) &= (K * v)(t) \\
&= \int_\R K(t-t') v(t') \dd t'
\end{aligned}
\end{equation}

\end{definition}

\begin{remark}[Boundedness via Young's Inequality]
We must verify that $W_K$ actually maps $L^2$ to $L^2$. By \textit{Young's convolution inequality}, if $K \in L^1(\R)$ and $v \in L^2(\R)$, then
\[
\begin{gathered}
K * v \in L^2(\R) \\
\|K * v\|_{L^2} \leq \|K\|_{L^1} \|v\|_{L^2}.
\end{gathered}
\]

Thus, the condition $K \in L^1$ is sufficient to guarantee that $W_K$ is a bounded (continuous) operator on $L^2(\R)$.
\end{remark}

\begin{proposition}[Action on Complex Exponentials]
Consider the complex exponential functions $e_\omega(t) = e^{i\omega t}$ for $\omega \in \R$. The action of $W_K$ on $e_\omega$ yields:
\begin{equation}
(W_K e_\omega)(t) = \hat{K}(\omega) e_\omega(t)
\end{equation}
Here the Fourier transform of $K$ is defined by
\[
\hat{K}(\omega) = \int_\R K(\tau) e^{-i\omega \tau} \dd \tau
\]

\end{proposition}

\begin{proof}
We evaluate the integral explicitly using the substitution $\tau = t - t'$, which implies $t' = t - \tau$ and $\dd t' = -\dd \tau$. The limits of integration change from $t' \in (-\infty, \infty)$ to $\tau \in (\infty, -\infty)$:
\begin{align}
(W_K e_\omega)(t) &= \int_{-\infty}^\infty K(t-t') e^{i\omega t'} \dd t' \notag \\
    &= \int_{\infty}^{-\infty} K(\tau) e^{i\omega (t-\tau)} (-\dd \tau) \notag \\
    &= \int_{-\infty}^\infty K(\tau) e^{i\omega t} e^{-i\omega \tau} \dd \tau \notag \\
    &= e^{i\omega t} \left( \int_{-\infty}^\infty K(\tau) e^{-i\omega \tau} \dd \tau \right) \notag \\
    &= \hat{K}(\omega) e_\omega(t)
\end{align}
This shows that $e_\omega$ is an eigenfunction of the convolution operator with eigenvalue $\hat{K}(\omega)$.
\end{proof}

\begin{remark}[Why $e_\omega$ are not true Eigenvectors in $L^2$]
For $e_\omega$ to be a valid eigenvector in the Hilbert space $L^2(\R)$, it must belong to $L^2(\R)$. However, its norm is:
\begin{equation}
\begin{aligned}
\|e_\omega\|^2 &= \int_{-\infty}^\infty |e^{i\omega t}|^2 \dd t \\
&= \int_{-\infty}^\infty 1 \dd t \\
&= \infty
\end{aligned}
\end{equation}

Thus, $e_\omega \notin L^2(\R)$. They are ``generalized eigenvectors''. Because the eigenvalues $\hat{K}(\omega)$ form a continuous range (as $\omega$ varies continuously over $\R$), the operator $W_K$ possesses a \textit{continuous spectrum}.
\end{remark}

\begin{proposition}[Convolution Operators are Never Compact]
Unless $W_K$ is the zero operator, it is never compact.
\end{proposition}

\begin{proof}
We prove this by constructing a bounded sequence with no convergent subsequence in the image.
Let $v \in L^2(\R)$ be any nonzero function. Define a sequence of translated functions $v_n(t) = v(t - n)$ for $n \in \N$.
Because translation preserves the $L^2$ norm,
\[
\|v_n\|_{L^2} = \|v\|_{L^2},
\]
so $\{v_n\}$ is a bounded sequence.
The image of this sequence under $W_K$ is:
\begin{equation}
\begin{aligned}
(W_K v_n)(t) &= \int_\R K(t-t') v(t'-n) \dd t' \\
&= \int_\R K(t - n - \tau) v(\tau) \dd \tau \\
&= (W_K v)(t - n)
\end{aligned}
\end{equation}

Thus, $W_K v_n$ is simply the function $W_K v$ shifted by $n$.
For $j \neq k$, the squared distance between two terms in the image sequence is:
\begin{equation}
\|W_K v_j - W_K v_k\|^2 = \int_\R |(W_K v)(t-j) - (W_K v)(t-k)|^2 \dd t
\end{equation}

As $j, k \to \infty$, the overlap between the shifted functions $(W_K v)(t-j)$ and $(W_K v)(t-k)$ approaches zero because $W_K v \in L^2(\R)$ (its mass cannot be infinite). Therefore, the cross-terms in the squared integral vanish, leaving:
\begin{equation}
\lim_{j,k \to \infty} \|W_K v_j - W_K v_k\|^2 = \|W_K v\|^2 + \|W_K v\|^2 = 2\|W_K v\|^2 > 0
\end{equation}

Since the distance between terms does not approach zero, the sequence $\{W_K v_n\}$ contains no Cauchy subsequence, and therefore no convergent subsequence. This violates the definition of a compact operator.
\end{proof}

\begin{remark}[Diagonalization via the Fourier Transform]
Although $W_K$ lacks a discrete eigenbasis, it is still ``diagonalized'' by the Fourier transform. The \textit{Convolution Theorem} states that
\[
\widehat{K * v}(\omega) = \hat{K}(\omega) \hat{v}(\omega).
\]

In the Fourier domain, the operator $W_K$ acts on $\hat{v}$ by simple pointwise multiplication by the function $\hat{K}(\omega)$. This is the exact continuous analogue of a diagonal matrix multiplying a vector. Rigorously, this is formalized by the theory of \textit{Direct Integrals}, where the discrete sum
\[
\sum_n \lambda_n \langle v, \phi_n \rangle \phi_n
\]
is replaced by the continuous integral
\[
\int_\R \hat{K}(\omega) \hat{v}(\omega) e^{i\omega t} \frac{\dd \omega}{2\pi}.
\]

\end{remark}

\section{Sobolev Spaces}

\subsection{Unbounded Differential Operators}

In quantum mechanics and field theory, we frequently encounter differential operators, such as the momentum operator $P = -i(\dd/\dd t)$. However, differential operators are inherently \textit{unbounded} on $L^2(\R)$.

\begin{remark}[Why Differential Operators are Unbounded]
An operator $D$ is bounded if there exists a constant $C$ such that $\|Dv\| \leq C\|v\|$ for all $v$. Consider the sequence of functions
\[
v_n(t) = (1/\sqrt{n}) \sin(nt) \chi(t),
\]
where $\chi(t)$ is a smooth bump function with compact support.
The $L^2$ norm of $v_n$ scales as
\[
\|v_n\|_{L^2} \sim (1/\sqrt{n}) \to 0.
\]
However, the derivative is
\[
v_n'(t) = \sqrt{n} \cos(nt) \chi(t) + \dots,
\]
so its $L^2$ norm scales as
\[
\|v_n'\|_{L^2} \sim \sqrt{n} \to \infty.
\]
Because the ratio
\[
\|v_n'\| / \|v_n\| \sim n \to \infty,
\]
no finite constant $C$ can bound the derivative operator. An unbounded operator cannot be defined on the entire Hilbert space; it requires a restricted, dense domain.
\end{remark}

To rigorously define the domain of differential operators and ensure they are closed and self-adjoint, we introduce Sobolev spaces. These spaces restrict the domain to functions that possess a specific degree of ``smoothness'' or ``regularity''.

\subsection{Fourier Characterization of Sobolev Regularity}

Instead of defining smoothness via pointwise derivatives (which can be problematic for non-differentiable functions), we define it via the decay of the Fourier transform at high frequencies. Rapid decay in the frequency domain corresponds to smoothness in the time domain.

\begin{definition}[Sobolev Space $H^s(\R)$]
For any real number $s \geq 0$, the Sobolev space $H^s(\R)$ is defined as the set of all functions $v \in L^2(\R)$ whose Fourier transform $\hat{v}(\omega)$ satisfies:
\begin{equation}
\|v\|_{H^s}^2 := \int_\R (1 + \omega^2)^s |\hat{v}(\omega)|^2 \dd \omega < \infty
\end{equation}
This space is a Hilbert space equipped with the inner product:
\begin{equation}
\langle v, u \rangle_{H^s} = \int_\R (1 + \omega^2)^s \hat{v}(\omega) \overline{\hat{u}(\omega)} \dd \omega
\end{equation}

\end{definition}

\begin{remark}[Equivalence to Weak Derivatives for Integer $s$]
Why does this Fourier definition control derivatives? We use the property that the Fourier transform turns differentiation into algebraic multiplication. Specifically, if we define the Fourier transform as
\[
\hat{v}(\omega) = \frac{1}{\sqrt{2\pi}} \int_\R v(t) e^{-i\omega t} \dd t,
\]
then the Fourier transform of the derivative is
\[
\widehat{v'}(\omega) = i\omega \hat{v}(\omega).
\]
By the \textit{Plancherel Theorem}, the $L^2$ norm is preserved under the Fourier transform:
\[
\|\hat{f}\|_{L^2} = \|f\|_{L^2}.
\]
Let us evaluate the $H^1$ norm explicitly:
\begin{align}
\|v\|_{H^1}^2 &= \int_\R (1 + \omega^2) |\hat{v}(\omega)|^2 \dd \omega \notag \\
    &= \int_\R |\hat{v}(\omega)|^2 \dd \omega + \int_\R \omega^2 |\hat{v}(\omega)|^2 \dd \omega \notag \\
    &= \|\hat{v}\|_{L^2}^2 + \|i\omega \hat{v}\|_{L^2}^2 \notag \\
    &= \|v\|_{L^2}^2 + \|\widehat{v'}\|_{L^2}^2 \notag \\
    &= \|v\|_{L^2}^2 + \|v'\|_{L^2}^2
\end{align}

This explicitly proves that $v \in H^1(\R)$ if and only if both $v$ and its first derivative $v'$ are in $L^2(\R)$. For a general integer $s$, expanding $(1+\omega^2)^s$ yields a polynomial in $\omega^2$ of degree $s$, which corresponds exactly to controlling the $L^2$ norms of all weak derivatives up to order $s$.
\end{remark}

\begin{definition}[Weak Derivative]
Because functions in $H^s(\R)$ might not be classically differentiable everywhere, we rely on \textit{weak derivatives}. A locally integrable function $g$ is the weak derivative of $f$ if for every smooth test function $\varphi$ with compact support:
\begin{equation}
\int_\R g(t) \varphi(t) \dd t = - \int_\R f(t) \varphi'(t) \dd t
\end{equation}

The Fourier definition of $H^s(\R)$ automatically handles weak derivatives because it only requires the integral of $|\omega|^s |\hat{v}(\omega)|^2$ to be finite, bypassing the need for pointwise classical derivatives.
\end{definition}

\begin{remark}[Fractional Differentiability for Non-Integer $s$]
For non-integer $s$ (e.g., $s = 1/2$), the weight $(1+\omega^2)^s$ does not correspond to an integer number of classical derivatives. Instead, it defines \textit{fractional differentiability}.
The factor $(1+\omega^2)^s$ penalizes high frequencies smoothly rather than abruptly. This measures the ``fractional smoothness'' of the function, interpolating between integer derivative spaces. This is crucial in the analysis of partial differential equations and neural field models, where solutions often exhibit regularity that is not an integer number of derivatives, but still possesses a well-defined, quantifiable degree of smoothness that ensures the convergence of numerical schemes and the well-posedness of the equations.
\end{remark}
